\subsection{Milestone 8F: Integration of the dependency-feedback monitor}

\subsubsection*{Milestone 8F-A: Signal and qualification boundary}

I integrated the selected five-bit monitor beside the official RAW dependency
checkers in \texttt{issue\_read\_operands.sv}. The monitor reuses the RAW-valid
decision, the complete matching-entry vector, the selected producer index, the
scoreboard issue pointer, and the live-entry vector already computed by CVA6.

I qualified every observation so that the monitor accepts only a real source
operand from a valid instruction inside the selected payload. Flushed and
exception-bearing instructions are rejected. RS1 is rejected when it is
replaced by a Z-immediate, and RS2 is rejected when the final operand mux
replaces it with an ordinary immediate. A stalled valid consumer is still
observed because an unresolved dependency is part of the internal behavior I
want to describe.

This exposed and corrected a limitation in my earlier replay traces. They had
marked every RS2 checker invocation as a used operand. In the ten new full-core
traces, 9433 unused RS2 rows were identified and excluded from feedback.

\subsubsection*{Milestone 8F-B: RTL integration}

I added the parameterized monitor to:

\begin{verbatim}
core/dependency_feedback_monitor.sv
\end{verbatim}

The monitor receives two source-operand observations per CVA6 issue port.
Events occurring simultaneously are combined before updating one five-bit
sticky bitmap. The RAW checker now exports its existing internal match vector,
so I did not duplicate the register-comparison logic only for the monitor.

The integrated CVA6 design compiled successfully with the project-pinned
Verilator 5.008. The monitor does not feed any instruction, forwarding, stall,
or architectural-state signal. As a non-interference check, the retired-
instruction log of \texttt{independent} remained byte-for-byte identical to its
pre-integration golden log. Both logs contained 1556 lines.

\subsubsection*{Milestone 8F-C: Verilator feedback readout}

The Verilator run accepts the following options:

\begin{verbatim}
+dependency_feedback_file=<output.csv>
+dependency_feedback_pc_start=<payload start address>
+dependency_feedback_pc_end=<payload end address>
\end{verbatim}

The PC interval is half-open. I obtain its addresses from the ELF symbols
\texttt{dependency\_payload\_start} and
\texttt{dependency\_payload\_end}. The feedback file contains the initial zero
and only subsequent changes of the sticky bitmap. Its last row is therefore the
final summary of the generated payload.

I first tested the monitor without a payload window. All ten programs produced
\texttt{0x1f}. This happened because the common startup, trap, and exit code
saturated the small bitmap before the differences between payloads could be
used. This result showed that the measurement window is a necessary part of the
feedback definition.

After restricting the monitor to the generated payload, all ten programs
passed and the integrated results matched the independent Python model:

\begin{center}
\begin{tabular}{lccc}
\hline
Program & Integrated & Offline & Result \\
\hline
\texttt{independent}               & \texttt{0x00} & \texttt{0x00} & MATCH \\
\texttt{load\_use}                 & \texttt{0x19} & \texttt{0x19} & MATCH \\
\texttt{same\_cycle\_bypass}      & \texttt{0x00} & \texttt{0x00} & MATCH \\
\texttt{scoreboard\_forward}      & \texttt{0x09} & \texttt{0x09} & MATCH \\
\texttt{newest\_producer\_control}& \texttt{0x19} & \texttt{0x19} & MATCH \\
\texttt{older\_producer}          & \texttt{0x1d} & \texttt{0x1d} & MATCH \\
\texttt{single\_producer\_control}& \texttt{0x1d} & \texttt{0x1d} & MATCH \\
\texttt{multiple\_producers}      & \texttt{0x1f} & \texttt{0x1f} & MATCH \\
\texttt{wrap\_pair\_early}       & \texttt{0x0b} & \texttt{0x0b} & MATCH \\
\texttt{wrap\_pair\_wrapped}     & \texttt{0x05} & \texttt{0x05} & MATCH \\
\hline
\end{tabular}
\end{center}

The corrected \texttt{scoreboard\_forward} result is \texttt{0x09}, rather
than the previous \texttt{0x0f}, because unused RS2 fields no longer generate
false wrapped and non-wrapped events. The corpus still produces seven distinct
five-bit summaries, so this correction removes false information without
reducing the observed treatment separation.

\subsubsection*{Result}

Milestones 8F-A, 8F-B, and 8F-C show that the selected five-bit feedback can be
computed inside the real CVA6 dependency-resolution logic and exported through
the official Verilator flow without changing retired execution. They also show
that compact feedback must describe the generated payload, not the common
runtime surrounding it.

The next step is to measure the incremental synthesis cost of the integrated
multi-operand cone and then use this five-bit value in a controlled MicroGP
experiment against a software-only baseline.
